\section{引言}
\label{sec:introduction}

空气质量监测网络需要在目标传感器停止上报时继续预测污染物浓度。例如，某站点过去一天的 PM2.5 和 PM10 测量均不可用，但更早的历史、本地气象变量以及其他站点的同期读数仍可能存在。预训练预测器可以直接接收这段不完整历史，修复方法也可以先推断缺失目标值再执行预测。处理方式的价值取决于其能否利用预测起点之前可获得的信息，降低后续预测误差。

已有方法提供了相关处理能力。SAITS 和 TimeMixer++ 等神经填补器学习时间与多变量表示，MoTM 则迁移预训练填补表示~\citep{du2023saits,wang2025timemixerpp,motmsoftware}。填补质量与下游预测质量可能不一致~\citep{lemorvan2021good}，时序填补基准也已包含下游评价~\citep{du2024tsibench}。Chronos-2 能够接收多变量历史，并在关联序列之间共享信息~\citep{ansari2025chronos2}。在这些能力之上，实际部署仍需确定修复历史中的哪些位置、保留哪些原观测，以及如何将统计预测与冻结模型配合使用。

本文研究目标传感器缺测后的固定预测方法。输入包含用于拟合简单统计模型的观测历史前缀，以及截至当前预测起点、长度更长且可能不完整的历史。方法只修复近期目标缺口，保留可用的长历史，并融合冻结预测器与统计模型的点预测。优化关注下游 MAE，同时完整报告 MSE。预测骨干参数保持不变，当前隐藏测量与待预测未来均不参与方法计算。

上述设计需要处理两个问题。首先，推断缺失目标时可用的变量集合随时间变化。替换全部变量可能改变有用的辅助信息，而仅向预测器提供较短的修复窗口又会丢失更早的观测。其次，合理的修复轨迹仍可能诱发不准确的预测。独立统计预测利用相同的预测前信息形成另一种估计，但融合是否有效取决于两个分量的误差关系。因此，实验需要同时保留匹配的长历史对照和统计融合对照。

本文定义\method{}（Peer Repair and Fusion）。该方法利用前缀拟合的岭回归，根据当前可观测的本地变量与同伴序列估计目标缺失值，将其写入受保护的长历史，再将 Chronos-2 预测与前缀拟合的向量自回归（vector autoregression，VAR）预测等权融合。在上述 PM 缺测实例中，仍在上报的同伴序列用于估计近期不可用的 PM 值，更早的目标观测继续保留在模型上下文中。VAR 分支对近期原观测进行滤波并外推最终状态。两个分支均截至同一预测起点，融合权重在全部评价任务中保持固定。

本文的工作包括三个方面。第~\ref{sec:method}~节给出可执行的修复与融合过程，明确拟合、观测保持和可用信息边界。第~\ref{sec:properties}~节分析其输入保持性质及融合降低平方误差的标准条件。第~\ref{sec:experiments}~节在自然缺测、人工目标缺测和更长预测跨度下评价同一固定方法，并补充停车位可用量的组件迁移结果。在北京 12 个站点的 45 个自然缺测窗口上，该方法的 MAE/MSE 为 0.388053/0.244688，相对匹配的“长历史原生预测＋VAR”对照降低 5.98\%/7.26\%。人工缺测和 96 步预测条件下存在更强的固定对照，因而这些结果限定了当前结论的适用范围。方法及应用重点由开发证据选定，其解释范围在实验设计中明确说明。

\section{相关工作}
\label{sec:related}

\paragraph{填补与下游预测。}
\citet{lemorvan2021good} 分析了协变量缺失时填补与预测之间的关系。SAITS 利用自注意力完成填补~\citep{du2023saits}；TimeMixer++ 为包括填补在内的多种时序任务提供表示~\citep{wang2025timemixerpp}；MoTM 通过连续建模实现预训练时序填补~\citep{motmsoftware}。TSI-Bench 在多类缺失条件下评价填补方法，并包含下游任务~\citep{du2024tsibench}。本文固定预测骨干，分别控制修复区域、预测历史长度及预测器直接接收的辅助变量。

\paragraph{缺失感知预测与预训练预测。}
Chronos-2 通过组注意力在关联序列、变量和协变量之间交换信息~\citep{ansari2025chronos2}。S4M 将缺失表示融入需要训练的预测架构~\citep{peng2025s4m}。本文采用既有 Chronos-2 参数，仅拟合外部统计模型，并以低秩适配对照~\citep{hu2021lora}检验更新少量新增参数能否提供更强的比较。关联序列条件化与模型适配均是已有能力，本文不将其作为原创的基本操作。

\paragraph{统计修复与预测组合。}
岭回归用于稳定相关预测变量下的线性估计~\citep{hoerl1970ridge}。线性预测组合已有长期研究~\citep{bates1969combination}，平均预测的平方误差分解也已有明确结论~\citep{krogh1994ensembles}。本文利用这些组件，为冻结多变量预测器定义具体的目标修复过程。实验检验该固定组合能否在相同长历史、相同 VAR 分支和较强替代修复之外获得收益，同时保留表现接近或更好的对照结果。
